\section{The Stack 到 StarCoder：代码数据的端到端管线}

这一章把一般 filtering protocol 落到真实 code corpus。代码数据的特殊性来自语言差异、仓库复制、generated files、secrets、benchmark solutions、metadata 与结构化开发事件。每一步都会改变模型能学到的任务。

\subsection{Raw 到 training data：大幅缩小不等于浪费}

The Stack v1 从约 6.4TB source code 缩到 StarCoder 使用的约 800GB；v2 从约 32.1TB 缩到约 3.6TB。这个 reduction 汇总语言选择、quality rules、dedup、PII 与格式约束。

\lecturefigure{slide-44.jpg}{StarCoder/The Stack v1：6.4TB raw 到 800GB training data}{Loubna Ben Allal 官方 deck，第 44 页}

\lecturefigure{slide-45.jpg}{StarCoder2/The Stack v2：32.1TB raw 到 3.6TB training data}{Loubna Ben Allal 官方 deck，第 45 页}

\begin{warningbox}{过滤率不是 quality metric}
删除 90\% 可能清掉 duplicates，也可能抹掉 rare languages；保留 90\% 可能保持 diversity，也可能保留 spam。必须报告被删类型、source/language retention、token distribution 与 downstream effect。
\end{warningbox}

\subsection{Language selection 与 per-language thresholds}

v1 从 358 languages 中选择约 86 个，排除 configs 和不维护语言；v2 扩展到 600+。平均行长、文件大小、comment ratio 等统计不能共用一个阈值，因为不同语言 idiom 和 generated-code patterns 不同。

\lecturefigure{slide-46.jpg}{The Stack v1：语言选择、社区抽样与 per-language filtering}{Loubna Ben Allal 官方 deck，第 46 页}

\begin{knowledgebox}{人在环路的具体形式}
BigCode 社区按 extension 检查约 100 个样本，记录真实 code、generated/vendor data、异常长度和格式，再推导 heuristics。它不是让志愿者逐文件审批，而是用 domain knowledge 校准 scalable rules。
\end{knowledgebox}

\teachervoice{00:30:00--00:33:20，讲者强调同一 average-line-length threshold 不能用于所有语言；社区 inspection 提供了按语言设计阈值的依据。}

\subsection{Near-dedup：MinHash 与 LSH 为什么适合 TB-scale code}

Exact dedup 只删完全相同文件；near-dedup 要识别改名、格式变化或局部修改。将文件表示为 shingles 集合 $A,B$，Jaccard similarity 为
\[
J(A,B)=\frac{|A\cap B|}{|A\cup B|}.
\]
MinHash 满足哈希碰撞概率近似 $\Pr[h(A)=h(B)]=J(A,B)$，再用 LSH（Locality-Sensitive Hashing）把相似 signatures 放进候选桶，避免全对全比较。

\lecturefigure{slide-47.jpg}{Near-deduplication：MinHash+LSH 与 StarCoder ablation 收益}{Loubna Ben Allal 官方 deck，第 47 页}

\begin{importantbox}{为什么 near-dedup 是高价值 filter}
它语言无关、容易并行，并直接减少重复训练与 memorization。slide 的 ablation 显示多种 code subsets 上 pass@1 明显提高，因此团队采用 strong dedup；这比未经验证的 stars/comments heuristic 更可靠。
\end{importantbox}

\teachervoice{00:33:20--00:35:20，讲者说 near-dedup 是效果最大的 filter，而且无需为 86 种语言分别修改；这同时满足收益与可扩展性。}

\subsection{PII、secrets 与 benchmark decontamination}

GitHub 自带 secret scanning 仍不能保证 corpus 无 emails、names、passwords 和 keys。BigCode 先人工标注 PII spans，再训练 StarPII NER model 扫描全量数据；课堂报告这一步约消耗 800 GPU hours。随后还要移除 evaluation benchmarks，避免 inflated score。

\lecturefigure{slide-48.jpg}{The Stack：PII removal 与 benchmark decontamination}{Loubna Ben Allal 官方 deck，第 48 页}

\begin{knowledgebox}{首次定义：PII 与 decontamination}
PII（Personally Identifiable Information）包括能识别个人的 names、emails、phones 等；secrets 还包括 tokens/passwords/keys。\textbf{Decontamination} 是从训练数据中移除 benchmark items 及其近似/派生版本，保护 evaluation independence。
\end{knowledgebox}

\begin{warningbox}{Detector 只能降低风险，不能出具“零 PII”证明}
NER 有 false negatives，代码字符串可能被混淆，历史 commit 可能泄漏删除内容。发布时还需 sample audits、reporting channel、gating 与 incident response。
\end{warningbox}

\teachervoice{00:34:00--00:36:20，讲者把 PII 处理的 annotation、NER 和 800 GPU-hour 成本说得很具体，也提醒 evaluation set 必须从 training corpus 移除。}

\subsection{Formatting：把 repository context 变成训练信号}

过滤后仍需决定 sequence format。StarCoder 在 code 前加入 repository/name/file metadata tokens；issues 包含 title、comment、user roles。格式让模型学到上下文边界，也会定义推理时可控制的接口。

\lecturefigure{slide-49.jpg}{StarCoder data formatting：repository、file、stars 与 issue metadata}{Loubna Ben Allal 官方 deck，第 49 页}

\lecturefigure{slide-50.jpg}{Git commits 与 Jupyter notebooks 的结构化格式}{Loubna Ben Allal 官方 deck，第 50 页}

对于 fill-in-the-middle（FIM），将文档切成 prefix $P$、middle $M$、suffix $S$，训练序列可重排为
\[
\langle\text{PRE}\rangle P\langle\text{SUF}\rangle S
\langle\text{MID}\rangle M,
\]
仍用 next-token loss，却让模型学习在已知前后文时生成中间 code。

\begin{warningbox}{Metadata tokens 也可能成为 shortcut}
若 stars、repository name 或 user identity 与质量标签高度相关，模型可能记忆 source bias。应评估移除 metadata 后的能力、privacy risk 和跨仓库 generalization。
\end{warningbox}

\subsection{StarCoder2 mixture：不同 source 贡献不同任务}

训练数据表列出 The Stack、GitHub issues、pull requests、notebooks、documentation、math/code datasets 等，并区分是否做 FIM、epochs 和 sample weight。读表时不要只看 token 百分比，要问每类 source 教什么 behavior。

\lecturefigure{slide-51.jpg}{StarCoder2 training-data mixture 与 source-specific settings}{Loubna Ben Allal 官方 deck，第 51 页}

若 source $i$ 有 $T_i$ tokens、sampling weight $w_i$，实际抽样概率为
\[
p_i=\frac{w_iT_i}{\sum_j w_jT_j}.
\]
小 source 可通过 $w_i>1$ 被 oversample，但重复 exposure 也增加 memorization，必须用 validation/evaluation 校准。

\begin{knowledgebox}{Source-to-behavior mapping}
Code files 教 syntax/algorithms；issues 教 problem description 与 discussion；commits/PRs 教 edit rationale；notebooks 教 code--output--markdown 交织；documentation 教 API usage；math datasets 教 symbolic/reasoning。Mixture 是能力设计，不只是拼接。
\end{knowledgebox}

\subsection{Tooling：可复现数据管线不能藏在一次性 notebook}

最后一页列出 datasets、datatrove、datatrove-like distributed pipelines、tokenizers/training code 与 BigCode processing repo。TB-scale filtering 需要 streaming、sharding（把文件或样本切成可由不同 workers 独立处理的分片）、checkpointing、deterministic transforms 和 dataset versioning。

\lecturefigure{slide-52.jpg}{数据过滤与训练 tooling：datasets、datatrove、训练框架与 BigCode code}{Loubna Ben Allal 官方 deck，第 52 页}

\begin{lstlisting}[caption={可复现数据管线的阶段契约}]
discover_sources_and_record_provenance()
normalize_and_language_classify()
apply_quality_filters_with_versioned_thresholds()
exact_and_near_deduplicate()
detect_pii_secrets_and_benchmark_overlap()
format_sequences_and_sample_mixture()
publish_manifest_hashes_stats_and_opt_out_state()
\end{lstlisting}

\subsection{本章小结}

StarCoder 的 data advantage 不是某个神奇 filter，而是 source selection、per-language inspection、near-dedup、PII/decontamination、structured formatting、mixture 和 reproducible tooling 的组合。每一步都有成本、失误方式和 evidence。

